\begin{frontmatter}

% TITLE PROPOSAL (per paper/DECISIONS.md: method = "Degree-Budget Topology Join",
% system = "CircuitDigitizer"). Adjust wording as preferred.
\title{CircuitDigitizer: Recovering Simulatable Netlists from Hand-Drawn Circuit Schematics with a Degree-Budget Topology Join}

\author[scaai]{Bosco Chanam\corref{cor1}}
\ead{bosco.chanam.btech2021@sitpune.edu.in}

\author[scaai]{Chris Dcosta}

\author[usc]{Pranavesh Kumar Talupuri}
\ead{talupuri@usc.edu}

\cortext[cor1]{Corresponding author}

\address[scaai]{Symbiosis Centre for Applied Artificial Intelligence, Department of Computer Science and Engineering, Symbiosis Institute of Technology, Pune, India}
\address[usc]{Thomas Lord Department of Computer Science, University of Southern California, Los Angeles, California, United States}

\begin{abstract}
Converting a photographed hand-drawn circuit schematic into a netlist that an
analog simulator can actually run is hard for a reason that is easy to overlook:
even when components and wire strokes are detected well, the recovered circuit is
almost always \emph{fragmented}. Faint or occluded strokes leave component
terminals unconnected, so the netlist breaks into electrically isolated pieces and
cannot be simulated. CircuitDigitizer is a transparent, stage-by-stage pipeline
that addresses this end to end. A trained oriented-bounding-box detector
(YOLO26M-OBB, 16 classes, 88.5\% mAP50) localizes components and guides occlusion
and region-of-interest cropping; a deterministic Sauvola-based computer-vision
stage extracts wire segments (point-based wire F1 = 0.9755 on 134 manually
verified images); and the central contribution, the \emph{Degree-Budget Topology
Join}, converts detected wires and components into a connected netlist. The join
first builds an endpoint graph and then treats every still-floating component
terminal --- the signature of a wire the detector dropped --- as a demand to be
satisfied by a reach-bounded minimum-cost matching, capped at one new edge per
terminal and guarded against self-loops. This recovers dropped connections while
the bounded edge budget structurally prevents the over-merging that shorts dense
layouts. Circuit connectivity rises from 68.8\% (a strong endpoint-graph baseline)
to 81.9\% with no per-image regressions, and the resulting netlists are emitted as
SPICE for DC operating-point simulation. Because real schematics have no net-level
ground truth, the join is also evaluated on a synthetic harness that authors known
circuits, lays them out as detector-style coordinate maps, injects detector-style
error, and scores the real join and simulator against the authored truth; the same
harness was used to discover the method. As a baseline, a strong vision-language
model (Mimo~2.5) prompted to read the same schematics end to end missed 38--94\% of
the topology, underscoring that explicit, inspectable topology reasoning remains
necessary. No wire-specific training is required for the extraction or join stages.
\end{abstract}

\begin{keyword}
Circuit diagram digitization \sep Netlist extraction \sep Topology join \sep
Hand-drawn schematics \sep SPICE simulation \sep Oriented object detection \sep
Document image analysis
\end{keyword}

\end{frontmatter}
